\subsection{Continuous finite-dimensional representations}

Let \(\varrho\) be a continuous linear representation of a topological group G in a finite-dimensional separated topological K-vector space E. We equip \(\mathcal{L}(\mathrm{E})\) with its unique separated topological vector space structure over K; the morphism \(\varrho\colon\mathrm{G}\to\mathrm{GL}(\mathrm{E})\) is then continuous since the topology of \(\mathrm{GL}(\mathrm{E})\) is induced by the topology of \(\mathcal{L}(\mathrm{E})\) , which coincides with the topology of simple convergence (EVT, I, p. 14, th. 2). Consequently, the representation \(\varrho\) is continuous. If G is a real Lie group, then the character of \(\varrho\) is an analytic function on G (LIE, III, p. 225, § 8, no. 1, th. 1). The contragredient representation is also continuous when E' is equipped with its unique topology as a separated topological vector space over K. Moreover, for every integer \(n\geqslant0\) the representations \(\mathrm{T}^{n}(\varrho)\) , \(\mathrm{S}^{n}(\varrho)\) and \(\wedge^{n}(\varrho)\) (loc. cit.) are continuous, when the corresponding spaces are equipped with their topologies as separated topological vector spaces over K.

Let \(\varrho_{1}\) and \(\varrho_{2}\) be continuous representations of a topological group G in separated topological vector spaces E \(_{1}\) and E \(_{2}\) of finite dimension. The linear representation \(\varrho_{1} \otimes \varrho_{2}\) of G in E \(_{1} \otimes\) E \(_{2}\) (LIE, III, p. 256, Appendix) is continuous, the space E \(_{1} \otimes\) E \(_{2}\) being equipped with its topology as a separated topological vector space over K.

\subsection{Irreducible Representations}

In this issue, G is a topological group.

DEFINITION 2. --- A representation \(\varrho\) of G in a topological K-vector space E is said to be irreducible if \(\{0\}\) is not dense in E and if the only subrepresentations of \(\varrho\) are \(\varrho\) and the representation in the closure of \(\{0\}\) .

If \(\varrho\) is an irreducible representation of G in a separated topological vector space E, then every nonzero element of E is a cyclic vector for \(\varrho\) .

Lemma 2. --- Let \(\pi\) and \(\varrho\) be continuous linear representations of G in separated topological K-vector spaces. We assume that \(\pi\) is irreducible. Every nonzero G-morphism from \(\pi\) to \(\varrho\) is injective and every nonzero G-morphism from \(\varrho\) to \(\pi\) has dense image.

In particular, if \(\pi\) and \(\varrho\) are irreducible and of finite dimension, then every nonzero G-morphism from \(\pi\) to \(\varrho\) is an isomorphism.

Since the space of \(\varrho\) is separated, the kernel of a G-morphism \(u\) from \(\pi\) to \(\varrho\) is closed, and hence defines a subrepresentation of \(\pi\) . If the morphism \(u\) is not zero, its kernel must therefore be reduced to \(\{0\}\) , since \(\pi\) is an irreducible representation in a separated topological vector space. Similarly, the closure of the image of a nonzero G-morphism from \(\varrho\) to \(\pi\) is a nonzero subrepresentation of \(\pi\) , hence is equal to the space of \(\pi\) .

The last assertion follows from what precedes.

Lemma 3. --- Let \(\varrho\) be a continuous linear representation of G in a separated topological K-vector space E of finite dimension. If E is nonzero, then there exists an irreducible subrepresentation of \(\varrho\) .

Since E is of finite dimension and nonzero, there exists a nonzero G-invariant subspace F of E of minimal dimension. This subspace is closed in E and defines a subrepresentation of E; every subrepresentation of F is also a subrepresentation of E, and the representation F is therefore irreducible by minimality.

Remark. --- A nonzero representation E of G does not always contain an irreducible subrepresentation (cf. V, p. 426, remark). One will see, however, that this is the case if G is compact, if K = C, and if E is a separated quasi-complete locally convex space over K and nonzero (prop. 7 of V, p. 464).

\subsection{Unitary representations}

DEFINITION 3. --- Let G be a topological group and E a Hilbert space over K. A unitary representation of G in E is a continuous linear representation \(\varrho\) of G in E such that, for every g in G, the endomorphism \(\varrho(g)\) of E is a unitary endomorphism (EVT, V, p. 40) of E.

In other words, a unitary representation is an isometric representation in a Hilbert space. In particular, a unitary representation is bounded.

The trivial representation of a topological group in a Hilbert space is unitary. Every subrepresentation of a unitary representation is unitary. The restriction to a subgroup of a unitary representation is unitary.

DEFINITION 4. --- Let G be a topological group and let \(\varrho\) and \(\pi\) be unitary representations of G in Hilbert spaces E and F respectively. A G-invariant sesquilinear form on E × F is a continuous sesquilinear form q on E × F such that

\[
q (x, y) = q (\varrho (g) x, \pi (g) y)
\]

for all \(g \in \mathrm{G}\) and every \((x, y) \in \mathrm{E} \times \mathrm{F}\) .

The vector space of G-invariant sesquilinear forms on E × F is denoted \(\mathrm{Sesq}_{\mathrm{G}}(\varrho,\pi)\) or \(\mathrm{Sesq}_{\mathrm{G}}(\mathrm{E},\mathrm{F})\) .

Example. --- Let G be a topological group and let \(\varrho\) be a unitary representation of G in E. The scalar product \(q\) on E is a G-invariant sesquilinear form on \(\mathrm{E} \times \mathrm{E}\) .

Lemma 4. --- Let G be a topological group and let \(\varrho\) be a homomorphism of G into the unitary group U(E) of a Hilbert space E. Then \(\varrho\) is a unitary representation of G if and only if \(\varrho\) is continuous at the element e of G for the topology of simple convergence on U(E). It suffices that this property hold for every \(x\) in a total subset of E.

The condition is evidently necessary. It is sufficient by remark 2 of INT, VIII, p. 129, § 2, no. 1, since the image of \(\varrho\) is equicontinuous in \(\mathcal{L}(\mathrm{E})\) and since the continuity of the map \(g \mapsto \varrho(g)x\) at \(e\) implies its continuity on G.

Let G be a topological group. If \(\varrho_{1}\) and \(\varrho_{2}\) are unitary representations of G and if u belongs to \(\operatorname{Hom}_{\mathrm{G}}(\varrho_{1},\varrho_{2})\) , then \(u^{*}\) belongs to \(\operatorname{Hom}_{\mathrm{G}}(\varrho_{2},\varrho_{1})\) . Indeed, since \(\varrho_{1}\) and \(\varrho_{2}\) are unitary representations, one has for every \(g\in G\)

\[
\begin{array}{r} u ^ {*} \circ \varrho_ {2} (g) = u ^ {*} \circ \varrho_ {2} (g ^ {- 1}) ^ {*} = (\varrho_ {2} (g ^ {- 1}) \circ u) ^ {*} \\ = (u \circ \varrho_ {1} (g ^ {- 1})) ^ {*} = \varrho_ {1} (g) \circ u ^ {*}. \end{array}
\]

Lemma 5. --- Let G be a topological group and let \(\varrho\) be a unitary representation of G in a complex Hilbert space E. The space \(\mathrm{Hom}_{\mathrm{G}}(\varrho, \varrho)\) is a unital stellar subalgebra of \(\mathcal{L}(\mathrm{E})\) .

The preceding shows that \(\operatorname{Hom}_{\mathrm{G}}(\varrho,\varrho)\) is a unital and self-adjoint subalgebra of \(\mathcal{L}(\mathrm{E})\) . Since it is closed in \(\mathcal{L}(\mathrm{E})\) , it is a stellar subalgebra of \(\mathcal{L}(\mathrm{E})\) .

Lemma 6. --- Let \(\pi\) and \(\varrho\) be unitary representations of a topological group G in Hilbert spaces E and F respectively. Let D be a dense subspace of E which is stable under \(\pi\) . Let u be a closed partial operator from E to F with domain D such that \(u \circ \pi(g) = \varrho(g) \circ u\) for all \(g \in G\) . Then the domain of \(u^{*}\) is stable under \(\varrho\) and one has the relation \(u^{*} \circ \varrho(g) = \pi(g) \circ u^{*}\) for all \(g \in G\) .

Let \(g \in \mathrm{G}\) and \(x \in \mathrm{dom}(u^*)\) . For every \(y \in \mathrm{dom}(u)\) , we have

\[
\begin{array}{r l} & {\langle \varrho (g) x \mid u (y) \rangle = \langle x \mid \varrho (g ^ {- 1}) u (y) \rangle = \langle x \mid u (\pi (g ^ {- 1}) y) \rangle} \\ & {\qquad = \langle u ^ {*} (x) \mid \pi (g ^ {- 1}) y \rangle = \langle \pi (g) (u ^ {*} (x)) \mid y \rangle ,} \end{array}
\]

since \(\varrho(g)^{*} = \varrho(g)^{-1} = \varrho(g^{-1})\) . This proves that \(\varrho(g)x \in \mathrm{dom}(u^{*})\) and that \(u^{*}(\varrho(g)x) = \pi(g)(u^{*}(x))\) . In particular, the domain of \(u^{*} \circ \varrho(g)\) contains the domain of \(u^{*}\) , and \(\pi(g) \circ u^{*} \subset u^{*} \circ \varrho(g)\) .

But moreover, if \(x \in \text{dom}(u^{*} \circ \varrho(g))\) , then \(x = \varrho(g^{-1})(\varrho(g)x)\) belongs to \(\text{dom}(u^{*})\) by virtue of the preceding applied to \(g^{-1}\) . We conclude that \(u^{*} \circ \varrho(g) = \pi(g) \circ u^{*}\) .

PROPOSITION 1. --- Let \(\varrho_{1}\) and \(\varrho_{2}\) be unitary representations of a topological group G in Hilbert spaces E \(_{1}\) and E \(_{2}\) . The map from Hom \(_{G}(\varrho_{1},\varrho_{2})\) to Sesq \(_{G}(\varrho_{2},\varrho_{1})\) which associates to u the sesquilinear form \(q_{u}\) defined by \(q_{u}(x,y)=\langle x|u(y)\rangle\) is an isomorphism of vector spaces.

If \(u\) is a G-morphism, then one has

\[
\begin{array}{c} q _ {u} (\varrho_ {2} (g) x, \varrho_ {1} (g) y) = \langle \varrho_ {2} (g) x \mid u (\varrho_ {1} (g)y) \rangle \\ = \langle \varrho_ {2} (g) x \mid \varrho_ {2} (g) u (y) \rangle = \langle x \mid u (y) \rangle = q _ {u} (x, y) \end{array}
\]

for all \(g \in G\) and every \((x, y) \in \mathrm{E}_{2} \times \mathrm{E}_{1}\) , hence the indicated map is a linear map from \(\operatorname{Hom}_{\mathrm{G}}(\varrho_{1}, \varrho_{2})\) to \(\operatorname{Sesq}_{\mathrm{G}}(\varrho_{2}, \varrho_{1})\) . By EVT, V, p. 16, cor. 2, it is injective.

Conversely, let \(q\) be a G-invariant sesquilinear form and \(u\) the unique linear map from \(\mathrm{E}_1\) to \(\mathrm{E}_2\) such that \(q(x,y) = \langle x \mid u(y) \rangle\) for all \((x,y) \in \mathrm{E}_2 \times \mathrm{E}_1\) (loc. cit.). For every \(g\) in \(G\) and every \((x,y)\)

in \(\mathrm{E}_2\times \mathrm{E}_1\) , we have

\[
\begin{array}{c} \langle x | (\varrho_ {2} (g) \circ u \circ \varrho_ {1} (g ^ {- 1})) (y) \rangle = \langle \varrho_ {2} (g ^ {- 1}) x | u (\varrho_ {1} (g ^ {- 1}) y) \rangle \\ = q (\varrho_ {2} (g ^ {- 1}) x, \varrho_ {1} (g ^ {- 1}) y) = q (x, y) = \langle x | u (y) \rangle , \end{array}
\]

hence \(u = \varrho_{2}(g) \circ u \circ \varrho_{1}(g^{-1})\) . Consequently, \(u \in \operatorname{Hom}_{\mathrm{G}}(\varrho_{1}, \varrho_{2})\) . The proposition follows.

Let \(\varrho\) be a unitary representation of a topological group G in a Hilbert space E. If K = C, denote by \(\overline{\mathrm{E}}\) the conjugate space of E (EVT, V, p. 6). If K = R, set \(\overline{\mathrm{E}} = \mathrm{E}\) . The conjugate representation \(\overline{\varrho}\) is the representation of G in \(\overline{\mathrm{E}}\) defined by \(\overline{\varrho}(g) = \varrho(g)\) for all \(g \in \mathrm{G}\) . It is a unitary representation of G; by definition, a subspace of \(\overline{\mathrm{E}}\) is a subrepresentation of \(\overline{\mathrm{E}}\) if and only if it is a subrepresentation of E.

PROPOSITION 2. --- Let \(\varrho\) be a unitary representation of a topological group G in a Hilbert space E. Denote by u the isometric isomorphism of \(\overline{\mathrm{E}}\) into \(\mathrm{E}'\) which associates to \(x\) the linear form \(y\mapsto \langle x|y\rangle\) .

\begin{enumerate}
\def\labelenumi{\alph{enumi})}
\item
  The map u is an isomorphism of the conjugate representation \(\overline{\varrho}\) into the contragredient representation \(\breve{\varrho}\) ;
\item
  Endow E' with the Hilbert space structure obtained by transporting structure via u; the contragredient representation \(\breve{\varrho}\) is a unitary representation of G in E'.
\end{enumerate}

By EVT, V, p. 15, th. 3 and the following remark, the map \(u\) is an isometric isomorphism.

For every g in G, every x in \(\overline{E}\) and every y in E, we have

\[
\begin{array}{c} \langle y, (\breve {\varrho} (g) \circ u) (x) \rangle = \langle \varrho (g ^ {- 1}) y, u (x) \rangle \\ = \langle x | \varrho (g ^ {- 1}) y \rangle = \langle \varrho (g) x | y \rangle = \langle y, u (\overline {{\varrho}} (g) x) \rangle , \end{array}
\]

hence \(\breve{\varrho}(g) \circ u = u \circ \overline{\varrho}(g)\) , which proves \(a\) ; the assertion \(b\) ) follows immediately.

COROLLARY. --- Let \(\varrho\) be a unitary representation of a topological group G in a Hilbert space E. The following conditions are equivalent:

\begin{enumerate}
\def\labelenumi{(\roman{enumi})}
\item
  The representation \(\varrho\) is irreducible;
\item
  The contragredient representation \(\breve{\varrho}\) is irreducible;
\item
  The conjugate representation \(\overline{\varrho}\) is irreducible.
\end{enumerate}

This follows from Proposition 2, and from the remark preceding it concerning subrepresentations of \(\overline{E}\) .

PROPOSITION 3. --- Let \(\varrho\) be a unitary representation of a topological group G in a Hilbert space E. Let \(\pi\) be a subrepresentation of \(\varrho\), with space F. The orthogonal F° of F in E is a subrepresentation of E such that E = F ⊕ F° and F° is isomorphic to E/F.

The space F° is closed. For every \(x \in F^{\circ}\) , all \(g \in G\) and every \(y \in F\) , we have \(\langle \varrho(g)x \mid y \rangle = \langle x \mid \varrho(g^{-1})y \rangle = 0\) since \(\varrho\) is unitary and F is a subrepresentation. Therefore \(\varrho(g)x \in F^{\circ}\) .

The canonical projection of E onto E/F is a G-morphism, so its restriction to F° is a G-morphism from F° to E/F; by EVT, V, p. 13, it is an isometric isomorphism.

We shall also denote by \(\pi^{\circ}\) the representation of G in F°.

PROPOSITION 4. --- Let \(\varrho\) be a unitary representation of a topological group G in a Hilbert space E. A closed subspace F of E is a subrepresentation of \(\varrho\) if and only if the orthogonal projector p of E with image F is a G-morphism from \(\varrho\) to \(\varrho\) .

Suppose that F is a subrepresentation of \(\varrho\) . Let \(x \in E\) and let us write \(x = p(x) + y\) where \(y \in F^{\circ}\) . For every \(g\) in G, we have

\[
\varrho (g) x = \varrho (g) (p (x)) + \varrho (g) y,
\]

and since \(\varrho(g)(p(x))\) belongs to F and \(\varrho(g)y\) to F° (prop. 3), we have \(p(\varrho(g)x) = \varrho(g)(p(x))\) . Hence \(p\) belongs to \(\operatorname{Hom}_{\mathrm{G}}(\varrho, \varrho)\) .

Conversely, if \(p \in \operatorname{Hom}_{\mathrm{G}}(\varrho, \varrho)\) then \(1_{E} - p\) is a G-morphism, hence \(\mathrm{F} = \operatorname{Ker}(1_{\mathrm{E}} - p)\) is a subrepresentation of \(\varrho\) .

\subsection{Hilbertian direct sum and tensor product of unitary representations}

Let G be a topological group. Let \((\varrho_{i})_{i\in\mathrm{I}}\) be a family of unitary representations of G in Hilbert spaces \(E_{i}\) . Let E be the space of the external Hilbert sum of the \(E_{i}\) . For every g in G and for every \(x=(x_{i})_{i\in\mathrm{I}}\) in E, we have

\[
\sum_ {i} \| \varrho_ {i} (g) x _ {i} \| ^ {2} = \sum_ {i} \| x _ {i} \| ^ {2} = \| x \| ^ {2},
\]

which proves that \((\varrho_i(g)x_i)_{i\in \mathrm{I}}\) is in E and has the same norm as \(x\) . Thus the element \(\varrho (g)\colon (x_{i})_{i\in \mathrm{I}}\mapsto (\varrho_{i}(g)x_{i})_{i\in \mathrm{I}}\) is a unitary element of \(\mathcal{L}(\mathrm{E})\) .

Lemma 7. --- The map \(g \mapsto \varrho(g)\) is a unitary representation of G in E.

By lemma 4 of V, p. 380, it suffices to show that for every i in I and every x in \(E_{i}\) , the map \(g \mapsto \varrho_{i}(g)x\) is continuous at the identity element e of G. This map is the composition of the continuous map \(g \mapsto \varrho_{i}(g)x\) and of the canonical injection of \(E_{i}\) into E. It is therefore continuous.

The representation \(\varrho\) is called the Hilbert sum of the unitary representations \((\varrho_{i})_{i\in I}\) ; it is denoted by \(\varrho = \bigoplus \varrho_{i}\) .

Let G and H be topological groups. Let \(\varrho_{1}\) (resp. \(\varrho_{2}\) ) be a unitary representation of G (resp. H) in a Hilbert space E \(_{1}\) (resp. E \(_{2}\) ). Let E = E \(_{1}\) \(\widehat{\otimes}_{2}\) E \(_{2}\) be the Hilbert tensor product space of E \(_{1}\) and E \(_{2}\) (EVT, V, p. 28, def. 1). For \((g,h)\in\mathrm{G}\times\mathrm{H}\) , let \(\varrho(g,h)\) be the continuous endomorphism \(\varrho_{1}(g)\widehat{\otimes}_{2}\varrho_{2}(h)\) of E (TVS, V, p. 28); we will denote it simply \(\varrho_{1}(g)\otimes\varrho_{2}(h)\) whenever no ambiguity is to be feared.

Lemma 8. --- a) The map \(\varrho\colon(g,h)\mapsto\varrho(g,h)\) is a unitary representation of G × H in E;

\begin{enumerate}
\def\labelenumi{\alph{enumi})}
\setcounter{enumi}{1}
\item
  For any orthonormal basis \((e_i)_{i\in \mathrm{I}}\) of \(\mathrm{E}_1\) the mapping of the Hilbert sum \(\bigoplus_{i\in \mathrm{I}}\mathrm{E}_2\) to E defined by \((y_i)_{i\in \mathrm{I}}\mapsto \sum_{i\in \mathrm{I}}e_i\otimes y_i\) is an isometric H-isomorphism of the Hilbert sum \(\bigoplus_{i\in \mathrm{I}}\varrho_2\) in \(\operatorname {Res}_{\{e\} \times \mathrm{H}}^{\mathrm{G}\times \mathrm{H}}(\varrho)\) ;
\item
  For any orthonormal basis \((f_j)_{j\in \mathrm{J}}\) of \(\mathrm{E}_2\) the mapping of the Hilbert sum \(\bigoplus_{j\in \mathrm{J}}\mathrm{E}_1\) to E defined by \((x_{j})_{j\in \mathrm{J}}\mapsto \sum_{j\in \mathrm{J}}x_{j}\otimes f_{j}\) is an isometric G-isomorphism of the Hilbert sum \(\bigoplus_{j\in \mathrm{J}}\varrho_1\) in \(\operatorname {Res}_{\mathbf{G}\times \{e\}}^{\mathbf{G}\times \mathbf{H}}(\varrho)\) .
\end{enumerate}

The map \((g,h)\mapsto\varrho(g,h)\) is a homomorphism from \(\mathrm{G}\times\mathrm{H}\) to \(\mathbf{GL}(\mathrm{E})\) (cf. EVT, V, p. 28, no. 2). Let \((e_{i})_{i\in\mathrm{I}}\) be an orthonormal basis of \(\mathrm{E}_{1}\) . By EVT, V, p. 29, prop. 3 and cor. 2, the map

\[
u \colon (y _ {i}) \mapsto \sum_ {i \in I} e _ {i} \otimes y _ {i}
\]

is an isometry from the Hilbert sum \(F_{2} = \bigoplus_{i \in I} E_{2}\) onto E. Via this isometry, the representation \(\varrho_{\mathrm{H}} \colon h \mapsto \varrho(e, h)\) of H in E is identified with the direct sum representation of the representations \((\varrho_{2})_{i\in I}\) in \(F_{2}\) . In particular, it is a unitary representation of H in E (lemma 7). Likewise, the homomorphism \(\varrho_{G}: g \mapsto \varrho(g,e)\) is a unitary representation of G in E.

Let \((g,h)\in\mathrm{G}\times\mathrm{H}\) . We have \(\varrho(g,h)=\varrho_{\mathrm{G}}(g)\circ\varrho_{\mathrm{H}}(h)\) , hence \(\varrho(g,h)\) is unitary; moreover, \(\varrho_{\mathrm{G}}(g)\) and \(\varrho_{\mathrm{H}}(h)\) commute, so lemma 1 of V, p. 377 implies that \(\varrho\) is a unitary representation of \(G\times H\) .

Finally, the assertions concerning the restriction of \(\varrho\) to the subgroups \(\{e\} \times \mathrm{H}\) and \(\mathrm{G} \times \{e\}\) were obtained in the course of the preceding argument.

The representation \((g,h)\mapsto\varrho_{1}(g)\otimes\varrho_{2}(h)\) of G×H is called the external tensor product of the unitary representations \(\varrho_{1}\) and \(\varrho_{2}\) , and is denoted \(\varrho_{1}\boxtimes\varrho_{2}\) .

Let \(n \in \mathbf{N}\) and let \((\mathrm{G}_i)_{1 \leqslant i \leqslant n}\) be a finite family of topological groups. Let \(\varrho_i\) be a unitary representation of \(\mathrm{G}_i\) in a Hilbert space \(\mathrm{E}_i\) for \(1 \leqslant i \leqslant n\) . One defines similarly a representation

\[
\varrho_ {1} \boxtimes \dots \boxtimes \varrho_ {n}
\]

of \(G_{1} \times \cdots \times G_{n}\) in the Hilbert space \(E = E_{1} \widehat{\otimes}_{2} \cdots \widehat{\otimes}_{2} E_{n}\) (EVT, V, p. 27).

Suppose that \(G_{i} = G\) for \(1 \leqslant i \leqslant n\) . Let \(\Delta_{n}: G \to G^{n}\) be the homomorphism defined by \(g \mapsto (g, \ldots, g)\) for all \(g \in G\) . We denote by \(\varrho_{1} \otimes \cdots \otimes \varrho_{n}\) the unitary representation \((\varrho_{1} \boxtimes \cdots \boxtimes \varrho_{n}) \circ \Delta_{n}\) of G. This is called the tensor product of the unitary representations \(\varrho_{i}\) .

For every permutation \(\sigma\) of \(\{1,\ldots,n\}\) , the canonical isomorphism

\[
\mathrm{E} _ {1} \widehat {\otimes} _ {2} \dots \widehat {\otimes} _ {2} \mathrm{E} _ {n} \rightarrow \mathrm{E} _ {\sigma (1)} \widehat {\otimes} _ {2} \dots \widehat {\otimes} _ {2} \mathrm{E} _ {\sigma (n)}
\]

(EVT, V, p. 28) is an isometric isomorphism

\[
\varrho_ {1} \otimes \dots \otimes \varrho_ {n} \rightarrow \varrho_ {\sigma (1)} \otimes \dots \otimes \varrho_ {\sigma (n)}
\]

of representations of G.

If \(\varrho_{i}=\varrho\) for \(1\leqslant i\leqslant n\) , where \(\varrho\) is a unitary representation of G, we also denote by \(\varrho^{\otimes n}\) the tensor product representation of the \(\varrho_{i}\) , and we say that it is the \(n\)th tensor power of \(\varrho\) .

\subsection{Matrix coefficients}

DEFINITION 5. --- Let G be a topological group and let \(\varrho\) be a unitary representation of G in a Hilbert space E. Let \(x\) and \(y\) be elements of E. The function from G to K given by \(g \mapsto \langle x | \varrho(g)y \rangle\) is called a matrix coefficient of \(\varrho\) , or a representative function. If \(x = y\) , one says that it is a diagonal matrix coefficient. If \(\varrho\) is of finite dimension, we say that it is a finite-dimensional matrix coefficient.

The matrix coefficients of \(\varrho\) are continuous and bounded functions on G. We denote by \(\Upsilon(G)\) (resp. \(\Theta(G)\) ) the set of matrix coefficients of complex unitary representations (resp. complex finite-dimensional representations) of G.

PROPOSITION 5. — Let G be a topological group. The sets \(\Theta(\mathrm{G})\) and \(\Upsilon(\mathrm{G})\) are unital involutive subalgebras of \(\mathcal{C}_{\mathrm{b}}(\mathrm{G})\) .

The constant function 1 is a matrix coefficient of the trivial representation of G on C. Let \(\varrho\) be a unitary representation of G and let \(x\) and \(y\) be vectors in the space of \(\varrho\) . For every \(\lambda \in \mathbf{C}\) and every \(g \in \mathrm{G}\) , we have \(\lambda \langle x \mid \varrho(g)y \rangle = \langle x \mid \varrho(g)(\lambda y) \rangle\) . Moreover, we have

\[
\langle x \mid \varrho (g) y \rangle = \langle \varrho (g ^ {- 1}) x \mid y \rangle = \overline {{\langle y \mid \varrho (g ^ {- 1}) x \rangle}}
\]

for all \(g \in G\) . Consequently, the sets \(\Theta(G)\) and \(\Upsilon(G)\) are stable under multiplication by scalars and under conjugation.

Let \(\varrho_{1}\) and \(\varrho_{2}\) be unitary representations of G; let \((x_{1}, y_{1})\) be vectors in the space of \(\varrho_{1}\) and \((x_{2}, y_{2})\) be vectors in the space of \(\varrho_{2}\) . For every \(g \in G\) , we then have

\[
\begin{array}{c} \langle x _ {1}   |   \varrho_ {1} (g) y _ {1} \rangle + \langle x _ {2}   |   \varrho_ {2} (g) y _ {2} \rangle = \langle (x _ {1}, x _ {2})   |   (\varrho_ {1} \oplus \varrho_ {2}) (g) (y _ {1}, y _ {2}) \rangle , \\ \langle x _ {1}   |   \varrho_ {1} (g) y _ {1} \rangle \langle x _ {2}   |   \varrho_ {2} (g) y _ {2} \rangle = \langle x _ {1} \otimes x _ {2}   |   (\varrho_ {1} \otimes \varrho_ {2}) (g) (y _ {1} \otimes y _ {2}) \rangle , \end{array}
\]

which proves that \(\Theta(G)\) and \(\Upsilon(G)\) are stable under addition and multiplication. The proposition follows.

\subsection{Schur's lemma}

In this issue, Hilbert spaces are complex.

PROPOSITION 6 (Schur's Lemma). --- Let \(\varrho\) be a unitary representation of a topological group G in a nonzero Hilbert space E. Then \(\varrho\) is irreducible if and only if \(\mathrm{Hom}_{\mathrm{G}}(\varrho, \varrho)\) is equal to \(\mathbf{C} \cdot 1_{\mathrm{E}}\) .

The space \(\mathrm{Hom}_{\mathrm{G}}(\varrho, \varrho)\) is a unital stellar subalgebra of \(\mathcal{L}(\mathrm{E})\) (lemma 5 of V, p. 380). For every subrepresentation F of E, the orthogonal projector \(p\) with image F is an idempotent element of the stellar algebra \(\mathrm{Hom}_{\mathrm{G}}(\varrho, \varrho)\) (prop. 4 of V, p. 383). If it is equal to \(\mathbf{C} \cdot 1_{\mathrm{E}}\) , this means that \(p = 0\) or \(p = 1_{\mathrm{E}}\) , which means that F = 0 or F = E. Thus \(\varrho\) is irreducible.

Conversely, suppose \(\varrho\) is irreducible. Let \(u\) and \(v\) be commuting elements of the stellar algebra \(\mathrm{Hom}_{\mathrm{G}}(\varrho, \varrho)\) such that \(uv = 0\) . Suppose \(u\) is nonzero. The kernel F of \(u\) then defines a subrepresentation of \(\varrho\) different from E, hence F is reduced to \(\{0\}\) ; since F contains the image of \(v\) , we deduce that \(v = 0\) . By Proposition 10 of I, p. 113, we therefore have \(\mathrm{Hom}_{\mathrm{G}}(\varrho, \varrho) = \mathbf{C} \cdot 1_{\mathrm{E}}\) .

COROLLARY 1. --- Let π be an irreducible unitary representation of a topological group G in a Hilbert space E. Let u be a closed partial operator on E. Suppose that u has dense domain, that dom(u) is stable under π, and that u ∘ π(g) = π(g) ∘ u for all g ∈ G. Then dom(u) = E and u is a scalar multiple of the identity.

The partial operator \(u^{*} \circ u\) is a positive self-adjoint partial operator on E (prop. 12 of IV, p. 241), the same holds for \(v = 1_{\mathrm{E}} + u^{*} \circ u\) and the latter is injective since \(-1 \notin \operatorname{Sp}(u^{*} \circ u)\) (prop. 17 of IV, p. 248). We have \(u^{*} \circ \pi(g) = \pi(g) \circ u^{*}\) for all \(g \in G\) (lemma 6 of V, p. 381), whence \(v \circ \pi(g) = \pi(g) \circ v\) for all \(g \in G\) . Since \(v\) is injective, it follows that \(v^{-1} \circ \pi(g) = \pi(g) \circ v^{-1}\) for all \(g \in G\) . But \(v^{-1}\) belongs to \(\mathcal{L}(\mathrm{E})\) , hence by prop. 6, there exists \(\lambda \in \mathbf{C}\) such that \(v^{-1} = \lambda 1_{\mathrm{E}}\) . Necessarily, \(\lambda \neq 0\) , which implies that \(\mathrm{E} = \mathrm{Im}(v^{-1}) = \mathrm{dom}(v) \subset \mathrm{dom}(u)\) , whence \(\mathrm{dom}(u) = \mathrm{E}\) . Since \(u\) is closed, we have \(u \in \mathcal{L}(\mathrm{E})\) , therefore \(u \in \mathrm{Hom}_{\mathrm{G}}(\pi, \pi)\) and \(u\) is a homothety by Prop. 6.

COROLLARY 2. --- Let \(\varrho\) and \(\pi\) be irreducible unitary representations of a topological group G in Hilbert spaces E and F, respectively. The space \(\mathrm{Hom}_{\mathrm{G}}(\varrho, \pi)\) is of dimension 1 if \(\varrho\) is isomorphic to \(\pi\) and is zero otherwise. In particular, if \(\varrho\) is isomorphic to \(\pi\) , every nonzero G-morphism from \(\varrho\) to \(\pi\) is an isomorphism.

Suppose there exists a nonzero G-morphism \(u\) from \(\varrho\) to \(\pi\) . The linear map \(u^{*} \circ u\) is an element of \(\mathrm{Hom}_{\mathrm{G}}(\varrho, \varrho)\) ; therefore there exists a complex number \(\lambda\) such that \(u^{*} \circ u = \lambda \cdot 1_{\mathrm{E}}\) (prop. 6).

We then have

\[
\langle u (x) \mid u (y) \rangle = \langle x \mid u ^ {*} u (y) \rangle = \lambda \langle x \mid y \rangle
\]

for all \(x\) and \(y\) in E. In particular, \(\lambda \neq 0\) since \(u\) is nonzero. Since \(\| u(x)\| = |\lambda |^{1 / 2}\| x\|\) for all \(x\in \mathrm{E}\) , the linear map \(u\) is injective, and the image of \(u\) is closed in F (Lemma 8 of I, p. 107); it is then a nonzero subrepresentation of the irreducible representation \(\pi\) , whence we deduce that \(u\) is surjective. Thus, \(u\) is an isomorphism from \(\varrho\) to \(\pi\) . The map \(v\mapsto u^{*}\circ v\) is then an isomorphism from the space \(\mathrm{Hom}_{\mathrm{G}}(\varrho ,\pi)\) to the space \(\mathrm{Hom}_{\mathrm{G}}(\varrho ,\varrho)\) , which is of dimension 1 (prop. 6).

COROLLARY 3. --- Let \(\varrho\) be a unitary representation of a topological group G in a Hilbert space E. The representation \(\varrho\) is irreducible if and only if the space \(\mathrm{Sesq}_{\mathrm{G}}(\varrho, \varrho)\) has dimension 1. This space is then generated by the scalar product on E.

In view of prop. 1 of V, p. 381, this follows from prop. 6.

COROLLARY 4. --- Let \(\varrho_{1}\) and \(\varrho_{2}\) be non-isomorphic irreducible unitary representations of a topological group G in Hilbert spaces E \(_{1}\) and E \(_{2}\) . The space Sesq \(_{G}\) ( \(\varrho_{1}, \varrho_{2}\) ) is zero.

By prop. 1 of V, p. 381, this follows from cor. 2.

COROLLARY 5. --- Let \(\varrho\) be a unitary representation of a topological group G in a Hilbert space E. Let \(\pi\) be an irreducible unitary representation of G in a Hilbert space F.

\begin{enumerate}
\def\labelenumi{\alph{enumi})}
\item
  For every nonzero G-morphism u from \(\pi\) to \(\varrho\) , there exists \(\lambda \in \mathbf{R}_{+}^{*}\) such that \(\lambda u\) is isometric;
\item
  Every G-morphism u from π to ρ has closed image;
\item
  Every nonzero G-morphism v from \(\varrho\) to \(\pi\) is surjective.
\end{enumerate}

Let us prove assertion \(a\) ), which implies the assertion \(b\) ) (lemma 8 of I, p. 107). Since the morphism \(u\) is nonzero, it is injective (lemma 2 of V, p. 378). The formula \(q(x,y) = \langle u(x)|u(y)\rangle\) for \(x\) and \(y\) in F then defines a continuous sesquilinear form on F. Since \(u\) is a G-morphism, we have \(q(\pi(g)x,\pi(g)y) = q(x,y)\) for all \(g\) in G and \((x,y)\in\mathrm{F}\times\mathrm{F}\) , hence \(q\) is G-invariant. By Cor. 3, there exists \(\alpha\in\mathbf{R}_{+}^{*}\) such that

\[
q (x, y) = \langle u (x) | u (y) \rangle = \alpha \langle x | y \rangle
\]

for all \((x,y)\in\mathrm{F}\times\mathrm{F}\) , and hence \(\alpha^{-1/2}u\) is isometric.

Let us finally prove assertion \(c\) ). Since \(\pi\) is irreducible, the image of \(v\) is dense in F (Lemma 2 of V, p. 378), hence the adjoint \(v^*\) is injective (TVS, V, p. 41, prop. 4). By \(b\) ), the image H of \(v^*\) is a closed subspace of E; it is therefore a subrepresentation of \(\varrho\) , and \(v^*\), induced by passage to the subspaces, is a G-isomorphism from F to H. The restriction \(w\) of \(v\) to H defines a G-morphism from H to F, which is injective since its kernel is the intersection of H and \(\text{Ker}(v) = \text{H}^\circ\) (loc. cit.). The map \(w\) is therefore an isomorphism (Corollary 2); in particular, \(v\) is surjective.

COROLLARY 6. --- Let \(G_{1}\) and \(G_{2}\) be topological groups. Let \(\varrho_{1}\) be an irreducible unitary representation of \(G_{1}\) in a Hilbert space \(E_{1}\) and \(\varrho_{2}\) an irreducible unitary representation of \(G_{2}\) in a Hilbert space \(E_{2}\) . The external tensor product \(\varrho_{1} \boxtimes \varrho_{2}\) is an irreducible unitary representation of \(G_{1} \times G_{2}\) in \(E_{1} \widehat{\otimes}_{2} E_{2}\) .

The external tensor product \(\varrho = \varrho_{1} \otimes \varrho_{2}\) is a unitary representation of \(G = G_{1} \times G_{2}\) in the space \(E = E_{1} \widehat{\otimes}_{2} E_{2}\) by lemma 8 of V, p. 384.

Let \(q\) be a sesquilinear form \((\mathrm{G}_1 \times \mathrm{G}_2)\)-invariant on E. For every pair \((x_1, y_1) \in \mathrm{E}_1^2\) , the map \((x_2, y_2) \mapsto q(x_1 \otimes x_2, y_1 \otimes y_2)\) belongs to \(\operatorname{Sesq}_{\mathrm{G}_2}(\varrho_2, \varrho_2)\) . Denote by \(b(x_1, y_1)\) the unique complex number such that

\[
q (x _ {1} \otimes x _ {2}, y _ {1} \otimes y _ {2}) = b (x _ {1}, y _ {1}) \left\langle x _ {2} \mid y _ {2} \right\rangle
\]

for all \((x_{2},y_{2})\in\mathrm{E}_{2}^{2}\) (cor. 3). Let \(\varepsilon\in E_{2}\) be of norm 1, so that \(b(x_{1},y_{1})=q(x_{1}\otimes\varepsilon,y_{1}\otimes\varepsilon)\) for all \((x_{1},y_{1})\in\mathrm{E}_{1}^{2}\) . This formula implies that the map b is sesquilinear on \(E_{1}\) . Moreover

\[
\| b (x _ {1}, y _ {1}) \| \leqslant \| q \| \| x _ {1} \otimes \varepsilon \| \| y _ {1} \otimes \varepsilon \| = \| q \| \| x _ {1} \| \| y _ {1} \|
\]

for all \((x_{1},y_{1})\in\mathrm{E}_{1}^{2}\) (EVT, V, p. 26, formula (5)), hence b is continuous.

Let \(g \in G\) and \((x_{1}, y_{1}) \in \mathrm{E}_{1}^{2}\) . Since \(\varrho_{2}\) is unitary and q is invariant, it follows that

\[
b \left(\varrho_ {1} (g) x _ {1}, \varrho_ {1} (g) y _ {1}\right) = q \left(x _ {1} \otimes \varrho_ {2} \left(g ^ {- 1}\right) \varepsilon , y _ {1} \otimes \varrho_ {2} \left(g ^ {- 1}\right) \varepsilon\right) = b \left(x _ {1}, y _ {1}\right),
\]

so that \(b \in \mathrm{Sesq}_{\mathrm{G}_1}(\varrho_1, \varrho_1)\) . There therefore exists a unique \(\lambda \in \mathbf{C}\) such that \(b(x_1, y_1) = \lambda \langle x_1 | y_1 \rangle\) for all \((x_1, y_1) \in \mathrm{E}_1^2\) (cor. 3), that is,

\[
q (x _ {1} \otimes x _ {2}, y _ {1} \otimes y _ {2}) = \lambda \langle x _ {1} | y _ {1} \rangle \langle x _ {2} | y _ {2} \rangle
\]

for all \((x_{1},y_{1},x_{2},y_{2})\in\mathrm{E}_{1}^{2}\times\mathrm{E}_{2}^{2}\) . From this one deduces that the space \(\operatorname{Sesq}_{\mathrm{G}_{1}\times\mathrm{G}_{2}}(\varrho,\varrho)\) has dimension 1, which implies that the representation \(\varrho=\varrho_{1}\boxtimes\varrho_{2}\) is irreducible (loc. cit.).

Recall that U denotes the group of complex numbers of modulus 1.

COROLLARY 7. --- Let \(\pi\) be an irreducible unitary representation of a topological group G in a Hilbert space E. There exists a continuous homomorphism \(\chi\) of the center C of G in U such that \(\pi(z) = \chi(z) \cdot 1_{\mathrm{E}}\) for all \(z \in \mathrm{C}\) . In particular, if G is commutative, we have \(\dim(\mathrm{E}) = 1\) .

For \(z \in \mathrm{C}\) , the map \(\pi(z)\) belongs to \(\mathrm{Hom}_{\mathrm{G}}(\pi, \pi)\) ; it is therefore of the form \(\chi(z) \cdot 1_{\mathrm{E}}\) for some complex number \(\chi(z)\) (prop. 6). Since \(\pi(z)\) is a unitary map, we have \(|\chi(z)| = 1\) . Moreover, since \(\pi\) is a homomorphism, the map \(z \mapsto \chi(z)\) is a homomorphism. Fix \(v \neq 0\) in E; the map \(z \mapsto \chi(z)v = \pi(z)v\) from C to U is continuous, and therefore the homomorphism \(\chi\) is continuous.

Finally, if G is commutative, then C = G, so \(\pi(g) = \chi(g) \cdot 1_{\mathrm{E}}\) for all g in G; every subspace of dimension 1 of E is then a subrepresentation of \(\pi\) , and since \(\pi\) is irreducible, the space E must be of dimension 1.

DEFINITION 6. --- Let π be an irreducible unitary representation of a topological group G in a Hilbert space E. The homomorphism χ of the center C of G into U such that π(z) = χ(z) · \(1_{\mathrm{E}}\) for all z in C is called the central character of π.

Note. --- Let \(\pi\) (resp. \(\varrho\) ) be an irreducible unitary representation of a topological group G (resp. H) in a Hilbert space E (resp. F). Denote by \(\chi\) (resp. \(\eta\) ) the central character of \(\pi\) (resp. \(\varrho\) ). The central character of the representation \(\overline{\pi}\) is \(\overline{\chi}\) and the central character of the irreducible unitary representation \(\pi \boxtimes \varrho\) of G × H (cor. 6) is the character \(\chi \boxtimes \eta\) : \((g, h) \mapsto \chi(g)\eta(h)\) of G × H.

\subsection{Semisimplicity}

DEFINITION 7. --- Let \(\varrho\) be a unitary representation of a topological group G in a Hilbert space E. We say that \(\varrho\) is semi-simple if there exists a family \((\mathrm{F}_i)_{i\in \mathrm{I}}\) of irreducible subrepresentations of \(\varrho\) such that E is the Hilbert sum of the subspaces \(\mathrm{F}_i\) .

One sometimes also says that a semisimple unitary representation admits a discrete decomposition or is discretely decomposable.

If \((\varrho_{i})_{i\in\mathrm{I}}\) is a family of semisimple unitary representations of a topological group G, then the Hilbert sum of the representations \(\varrho_{i}\) is semisimple.

PROPOSITION 7. --- Let G be a topological group and let \(\varrho\) be a unitary representation of G in a complex Hilbert space E. The representation \(\varrho\) is semi-simple if and only if every nonzero subrepresentation of \(\varrho\) contains an irreducible subrepresentation.

Suppose that \(\varrho\) is semisimple. Let \((\mathrm{F}_i)_{i\in \mathrm{I}}\) be a family of irreducible subrepresentations of \(\varrho\) such that E is the Hilbert sum of the subspaces \(\mathrm{F}_i\) . Let F be a nonzero subrepresentation of E. There exists \(i\in \mathrm{I}\) such that the restriction to F of the orthogonal projector with image \(\mathrm{F}_i\) is not zero. This orthogonal projector then defines, by passing to subspaces, a nonzero G-morphism \(p_i\) from F into \(\mathrm{F}_i\) . By Schur's lemma, the G-morphism \(p_i\) is surjective (cor. 5 of V, p. 388). The orthogonal complement in F of the kernel of \(p_i\) is a subrepresentation of F isomorphic to \(\mathrm{F}_i\) (prop. 3 of V, p. 383), hence irreducible.

Let us prove the converse assertion. Let \(\mathcal{F}\) be the set of closed subspaces F of E that are stable under G such that the subrepresentation of \(\varrho\) in F is irreducible. Let \(\mathcal{O}\) be the set of subsets \(\mathcal{G}\) of \(\mathcal{F}\) such that the subspaces belonging to \(\mathcal{G}\) are pairwise orthogonal.

The set \(\mathcal{O}\) is of finite character (E, III, p. 34, def. 2). Let then \(\mathcal{G}\) be a maximal element of \(\mathcal{O}\) (E, III, p. 35, th. 1). Let E \(_1\) be the subspace of E that is the Hilbert sum of the irreducible subrepresentations F of \(\mathcal{G}\) . If one had E \(_1\) ≠ E, the orthogonal complement of E \(_1\) would not be zero, and the subrepresentation of G in E \(_1^{\circ}\) would by hypothesis contain an irreducible subrepresentation. Its space F would be orthogonal to the elements of \(\mathcal{G}\) , so that \(\mathcal{G} \cup \{F\} \in \mathcal{O}\) ; this contradicts the maximality of \(\mathcal{G}\) , so E = E \(_1\) , which concludes the proof.

COROLLARY 1. --- Let G be a topological group and let \(\varrho\) be a semisimple unitary representation of G in a complex Hilbert space E. Every subrepresentation of \(\varrho\) (resp. every quotient representation of \(\varrho\) ) is semisimple.

If \(\varrho_{1}\) is a subrepresentation of \(\varrho\) , then every nonzero subrepresentation of \(\varrho_{1}\) contains an irreducible subrepresentation (prop. 7), hence \(\varrho_{1}\) is semi-simple (loc. cit.).

By prop. 3 of V, p. 383, every quotient representation of \(\varrho\) is isomorphic to a subrepresentation of \(\varrho\) ; therefore it is semisimple.

COROLLARY 2. --- Every finite-dimensional unitary representation \(\varrho\) of a topological group G is semisimple.

This follows from prop. 7 and lemma 3 of V, p. 379.

COROLLARY 3. --- Let \(\varrho_{1}\) and \(\varrho_{2}\) be finite-dimensional unitary representations of a topological group G. The representations \(\varrho_{1}\) and \(\varrho_{2}\) are isomorphic if and only if \(\chi_{\varrho_{1}} = \chi_{\varrho_{2}}\) .

This follows from Corollary 2 and from A, VIII, p. 389, prop. 1, b).

\subsection{Classes of unitary representations}

Lemma 9. --- Let E be a Hilbert space over K. Let F ⊂ E be a dense vector subspace of E. Then the Hilbert dimension of E is less than or equal to the dimension of F.

If E has finite Hilbert dimension, this is equal to the dimension of E. Suppose that E has infinite Hilbert dimension. The subspace F is then of infinite dimension. Let B be an orthonormal basis of E and let B' be a basis of F. For every \(x \in B\) , there exists an element \(f(x) \in B'\) such that \(\langle x \mid f(x) \rangle \neq 0\) , since otherwise we would have \(x \in F^{\circ} = \{0\}\) . One thus defines a map \(f: B \to B'\) .

For every \(y \in B'\) , the set \(f^{-1}(y)\) is contained in the set of \(x \in B\) such that \(\langle x | y \rangle \neq 0\) and is therefore countable (EVT, V, p. 21, prop. 4). By E, III, p. 50, prop. 4, one obtains \(\text{Card(B)} = \text{Card}(f(B)) \leqslant \text{Card(B')}\) .

Denote by \(Is_{G}(\pi_{1},\pi_{2})\) the relation

“G is a topological group and \(\pi_{1}\) and \(\pi_{2}\) are

unitary representations of G in Hilbert spaces over K that are

isomorphic.” With respect to \(\pi_{1}\) and \(\pi_{2}\) , this is an equivalence relation. For any unitary representation \(\pi\) of G in a Hilbert space over K, we shall denote by \(\mathrm{cl}(\pi)\) the equivalence class of \(\pi\) (cf. E, II, p. 47); it is therefore a unitary representation of G isomorphic to \(\pi\) ; one says that \(\mathrm{cl}(\pi)\) is the class of \(\pi\) . Unitary representations \(\pi_{1}\) and \(\pi_{2}\) in Hilbert spaces over K are isomorphic if and only if \(\mathrm{cl}(\pi_{1}) = \mathrm{cl}(\pi_{2})\) .

Let G be a topological group. Let c be a cardinal. The relation

“λ is a class of unitary representation of G in

a complex Hilbert space of Hilbert dimension \(\leqslant\mathfrak{c}\)”

is collectivizing in \(\lambda\) (E, II, p. 3). Indeed, every Hilbert space over K of dimension \(\leqslant \mathfrak{c}\) is isometrically isomorphic to a Hilbert subspace of \(\ell^2 (\mathfrak{c})\) (EVT, V, p. 23, cor. 2), and the assertion then follows from E, II, p. 47.

Let \(\pi\) be an irreducible unitary representation of G in a Hilbert space E. Let \(x\) be a nonzero element of E. Since \(\pi\) is irreducible, the vector \(x\) is a cyclic vector of \(\pi\) , which implies that the Hilbert dimension of E is \(\leqslant\) Card(G) (Lemma 9, applied to the dense subspace generated by the elements \(\pi(g)x\) for \(g \in G\) ). The classes of irreducible unitary representations of G in a Hilbert space over K therefore belong to the set of classes of unitary representations of G in a Hilbert space over K of Hilbert dimension \(\leqslant\) Card(G); consequently, they form a set.

DEFINITION 8. --- We denote by \(\widehat{\mathbf{G}}\) the set of classes of irreducible unitary representations of \(\mathbf{G}\) in a complex Hilbert space. We say that \(\widehat{\mathbf{G}}\) is the unitary dual of \(\mathbf{G}\) .

For every irreducible unitary representation \(\pi\) of G in a complex Hilbert space, we therefore have \(\operatorname{cl}(\pi) \in \widehat{\mathbf{G}}\) .

Remarks. --- 1) Suppose that G is a commutative group. By Corollary 7 of V, p. 390, every irreducible unitary representation of G in a complex Hilbert space has dimension 1. If G is locally compact, the set \(\widehat{G}\) is identified with the set of unitary characters of G (def. 1 of II, p. 201), and the notation \(\widehat{G}\) is therefore compatible with that introduced in Def. 2 of II, p. 201.

\begin{enumerate}
\def\labelenumi{\arabic{enumi})}
\setcounter{enumi}{1}
\item
  If G is finite, the set \(\widehat{\mathbf{G}}\) is in bijection with the set of classes of \(\mathbf{C}[\mathrm{G}]\) -simple modules (A, VIII, p. 47), which is also denoted \(\widehat{\mathbf{G}}\) in A, VIII, p. 396.
\item
  For \(\pi \in \widehat{\mathbf{G}}\) , we shall identify \(\overline{\pi}\) with the class in \(\widehat{\mathbf{G}}\) of the conjugate representation of \(\pi\) .
\item
  If \(\pi\) is an irreducible unitary representation of G in a finite-dimensional complex Hilbert space, its character \(\chi_{\pi}\) depends only on the class of \(\pi\) . One can therefore speak of the set of characters of finite-dimensional complex irreducible unitary representations of G.
\end{enumerate}

\subsection{Isotypic components}

DEFINITION 9. --- Let G be a topological group and let \(\pi\) be a continuous irreducible representation of G. Let \(\varrho\) be a continuous representation of G in a separated locally convex space E. The \(\pi\)-isotypic component of \(\varrho\) is called the closure in E of the sum of the spaces of all subrepresentations of \(\varrho\) isomorphic to \(\pi\) . This subspace is denoted \(\mathrm{M}_{\pi}(\varrho)\).

The space \(\mathrm{M}_{\pi}(\varrho)\) is a closed subspace of E, which defines a subrepresentation of \(\varrho\) . This space depends only on the class of \(\pi\) in \(\widehat{G}\) .

PROPOSITION 8. — Let G be a topological group. Let \(\varrho\) be a unitary representation of G in a complex Hilbert space E.

\begin{enumerate}
\def\labelenumi{\alph{enumi})}
\item
  Let \(\pi_1\) and \(\pi_2\) be non-isomorphic irreducible unitary representations of G. The subspaces \(\mathrm{M}_{\pi_1}(\varrho)\) and \(\mathrm{M}_{\pi_2}(\varrho)\) are orthogonal;
\item
  Let \(\varrho'\) be a unitary representation of G and \(u\) a G-morphism from \(\varrho\) to \(\varrho'\) . For every irreducible unitary representation \(\pi\) of G, we have \(u(\mathrm{M}_{\pi}(\varrho)) \subset \mathrm{M}_{\pi}(\varrho')\) .
\end{enumerate}

Let \(E_{1}\) and \(E_{2}\) be subspaces of E defining subrepresentations isomorphic to \(\pi_{1}\) and \(\pi_{2}\) respectively. The orthogonal projector of E with image \(E_{2}\) defines by passing to subspaces an element of \(\mathrm{Hom}_{\mathrm{G}}(\pi_{1},\pi_{2})\) , which is zero (cor. 2 of V, p. 387), which proves that \(E_{2}\) is orthogonal to \(E_{1}\) . This proves assertion \(a\) ).

For assertion \(b\) ), note that \(\mathrm{M}_{\pi}(\varrho)\) is, by definition, the closure in E of the space generated by the elements \(x \in \mathrm{E}\) belonging to a closed subspace \(\mathrm{F} \subset \mathrm{E}\) stable under \(\varrho\) such that the subrepresentation \(\varrho_{\mathrm{F}}\) of \(\varrho\) in F is isomorphic to \(\pi\) . It therefore suffices to prove that, in this case, one has \(u(x) \in \mathrm{M}_{\pi}(\varrho')\) . Let \(\mathrm{H} = u(\mathrm{F})\) ; it is a closed subspace of the space of \(\varrho'\) (cor. 5 of V, p. 388) stable under \(\varrho'\) , and u, induced by passage to the subspaces, is a surjective G-morphism from F to H. If H is not zero, then this G-morphism is an isomorphism by Corollary 2 of V, p. 387. The representation of G on H is therefore isomorphic to \(\pi\) , whence it follows that \(u(x)\) belongs to \(\mathrm{M}_{\pi}(\varrho')\) .

For every vector space H and every family \((\mathrm{H}_{i})_{i\in\mathrm{I}}\) of subspaces of H, we say that the spaces \((\mathrm{H}_{i})_{i\in\mathrm{I}}\) are in direct sum if the family \((\mathrm{H}_{i})_{i\in\mathrm{I}}\) satisfies the equivalent conditions of Prop. 11 of A, II, p. 18.

PROPOSITION 9. --- Let G be a topological group and let \(\varrho\) be a unitary representation of G in a complex Hilbert space E. Let \(\pi\) be an irreducible unitary representation of G in a complex Hilbert space \(\mathrm{E}_{\pi}\) . Denote by \(v\) the linear map from \(\mathrm{Hom}_{\mathrm{G}}(\pi, \varrho) \otimes \mathrm{E}_{\pi}\) into E such that \(v(u \otimes x) = u(x)\) for all \((u, x) \in \mathrm{Hom}_{\mathrm{G}}(\pi, \varrho) \times \mathrm{E}_{\pi}\) .

The linear map \(v\) is injective and its image is the sum of the spaces of all subrepresentations of \(\varrho\) which are isomorphic to \(\pi\) . In particular, the image of \(v\) is dense in \(M_{\pi}(\varrho)\) .

Corollary 5 of V, p. 388 implies that the image of v is the sum of the spaces of all subrepresentations of \(\varrho\) which are isomorphic to \(\pi\) .

Let us prove that \(v\) is injective. Let \((u_i)_{i \in I}\) be a basis of the vector space Hom \(_G(\pi, \varrho)\) . For \(i \in I\) , denote by F \(_i\) the image of \(u_i\) and \(\widetilde{u}_i \colon E_\pi \to F_i\) the G-isomorphism deduced from \(u_i\) by passage to the subspaces.

Let us first prove, by induction on the cardinality of a finite subset J of I, that the subspaces \((\mathrm{F}_{i})_{i\in\mathrm{J}}\) are in direct sum.

The case where J is empty is immediate. Suppose that J is nonempty and that the required property holds for subsets of I of cardinality at most \(\text{Card(J)} - 1\) .

Let \(j\) be a fixed element of J. The induction hypothesis implies that the subspaces \(\mathrm{F}_{i}\) for \(i\in\mathrm{J}-\{j\}\) are in direct sum. Let \(\mathrm{F}^{\prime}\) be their sum; it now suffices to prove that \(\mathrm{F}_{j}\cap\mathrm{F}^{\prime}=\{0\}\) .

Suppose that the intersection of \(F_{j}\) and of \(F'\) is not reduced to 0; this intersection is a subrepresentation of \(F_{j}\) , and since the latter is irreducible, we therefore have \(F_{j} \cap F' = F_{j}\) , whence \(F_{j} \subset F'\) .

For \(i \in \mathrm{J} - \{j\}\) , denote by \(\operatorname{pr}_i: \mathrm{F}' \to \mathrm{F}_i\) the projection and by \(\iota_i: \mathrm{F}_i \to \mathrm{F}'\) the inclusion. We have a canonical isomorphism

\[
\operatorname{Hom} _ {\mathrm{G}} \left(\mathrm{F} _ {j}, \mathrm{F} ^ {\prime}\right)\rightarrow \bigoplus_ {i \in \mathrm{J} - \{j \}} \operatorname{Hom} _ {\mathrm{G}} \left(\mathrm{F} _ {j}, \mathrm{F} _ {i}\right)
\]

(formula (1), p. 377) such that \(u \colon \mathrm{F}_j \to \mathrm{F}'\) has as image the family \((\mathrm{pr}_i \circ u)_{i \in \mathrm{J} - \{j\}}\) (A, II, p. 13, cor. 1). Corollary 2 of V, p. 387 implies that \(\widetilde{u}_i \circ \widetilde{u}_j^{-1}\) , which is nonzero, is a basis of the space \(\mathrm{Hom}_{\mathrm{G}}(\mathrm{F}_j, \mathrm{F}_i)\) .

Therefore, the family of G-morphisms \((\iota_{i} \circ \widetilde{u}_{i} \circ \widetilde{u}_{j}^{-1})_{i \in \mathrm{J} - \{j\}}\) is a basis of \(\mathrm{Hom}_{\mathrm{G}}(\mathrm{F}_{j}, \mathrm{F}')\) .

Denote by \(\iota\) the inclusion of \(F_{j}\) in \(F'\) ; it is an element of \(\mathrm{Hom}_{\mathrm{G}}(\mathrm{F}_{j}, \mathrm{F}')\) , hence it is a linear combination of the G-morphisms \(\iota_{i} \circ \widetilde{u}_{i} \circ \widetilde{u}_{j}^{-1}\) for \(i \in \mathrm{J} - \{j\}\) . It follows that \(u_{j}\) is a linear combination of the maps \(u_{i}\) for \(i \neq j\) , which contradicts the linear independence of the family \((u_{i})_{i \in \mathrm{I}}\) . The assertion is therefore proved by induction.

Let us now prove that v is injective. Let w be an element of \(\mathrm{Hom}_{\mathrm{G}}(\pi,\varrho)\otimes\mathrm{E}_{\pi}\) . There exists a unique family \((x_{i})_{i\in\mathrm{I}}\) in \(E_{\pi}\) with finite support such that

\[
w = \sum_ {i \in \mathrm{I}} u _ {i} \otimes x _ {i}
\]

(A, II, p. 62, cor. 1) and we then have

\[
v (w) = \sum_ {i \in I} u _ {i} (x _ {i}).
\]

From what precedes, the condition \(v(w) = 0\) therefore implies that \(u_{i}(x_{i}) = 0\) for all \(i \in I\) , whence \(x_{i} = 0\) for all i since \(u_{i}\) is injective, and therefore w = 0.

PROPOSITION 10. --- Let G be a topological group. Let \(\pi\) be an irreducible unitary representation of G in a complex Hilbert space E \(_{\pi}\) and let \(\varrho\) be a unitary representation of G in a complex Hilbert space E.

There exist families \((\mathrm{E}_{i})_{i\in\mathrm{I}}\) of closed invariant subspaces of E such that the subrepresentation of \(\varrho\) in \(E_{i}\) is isomorphic to \(\pi\) for all \(i\in I\) and such that \(\mathrm{M}_{\pi}(\varrho)\) is the Hilbert sum of the spaces \(E_{i}\) . Moreover, the cardinality of I is independent of the family \((\mathrm{E}_{i})_{i\in\mathrm{I}}\) satisfying these properties.

Let us prove the existence of families satisfying the stated conditions. Let \(\mathcal{O}\) be the set of subsets C of \(\mathrm{Hom}_{\mathrm{G}}(\pi, \varrho)\setminus\{0\}\) such that the subspaces \(u(\mathrm{E}_{\pi})\) for \(u \in \mathrm{C}\) are pairwise orthogonal. The set \(\mathcal{O}\) is ordered by inclusion. It is nonempty because the empty set is a member, and it is of finite character (E, III, p. 34, def. 2) since C belongs to \(\mathcal{O}\) if and only if every two-element subset of C belongs to \(\mathcal{O}\) . By E, III, p. 35, th. 1, there exists a maximal element C of \(\mathcal{O}\) . Let F be the closure of the subspace generated by the spaces \(u(\mathrm{E}_{\pi})\) for \(u \in \mathrm{C}\) ; it is the Hilbert sum of the spaces \(u(\mathrm{E}_{\pi})\) for \(u \in C\) . We shall prove that \(\mathrm{F} = \mathrm{M}_{\pi}(\varrho)\) , which will prove that the family \((u(\mathrm{E}_{\pi}))_{u \in \mathrm{C}}\) has the required properties.

By definition, \(u(\mathrm{E}_{\pi}) \subset \mathrm{M}_{\pi}(\varrho)\) for all \(u \in \mathrm{C}\) , hence \(\mathrm{F}\) is contained in \(\mathrm{M}_{\pi}(\varrho)\) . To prove the reverse inclusion, it suffices to prove that if \(v\) is a G-morphism from \(\pi\) to \(\varrho\) , then its image is contained in \(\mathrm{F}\) . Let \(p\) be the orthogonal projector of \(\mathrm{E}\) with image \(\mathrm{F}^{\circ}\) ; it is a G-morphism, since \(\mathrm{F}\) is a subrepresentation of \(\mathrm{E}\) (prop. 4 of V, p. 383). The image of the G-morphism \(p \circ v\) is orthogonal to \(\mathrm{F}\) ; it is therefore zero (otherwise \(\mathrm{C} \cup \{p \circ v\} \in \mathcal{O}\) , which contradicts the maximality of \(\mathrm{C}\) ), and consequently the image of \(v\) is contained in \(\mathrm{F}\) .

Now let \((\mathrm{E}_{i})_{i\in\mathrm{I}}\) and \((\mathrm{F}_{j})_{j\in\mathrm{J}}\) be families of pairwise orthogonal closed invariant subspaces of E, such that the subrepresentation of G in \(E_{i}\) (resp. in \(F_{j}\) ) is isomorphic to \(\pi\) for all \(i\in I\) (resp. for every \(j\in J\) ), and such that \(\mathrm{M}_{\pi}(\varrho)\) is the Hilbert sum of the family \((\mathrm{E}_{i})_{i\in\mathrm{I}}\) and of the family \((\mathrm{F}_{j})_{j\in\mathrm{J}}\) .

If I is finite, then

\[
\dim \operatorname{Hom} _ {\mathrm{G}} (\pi , \mathrm{M} _ {\pi} (\varrho)) = \dim \operatorname{Hom} _ {\mathrm{G}} \Bigl (\mathrm{E} _ {\pi}, \bigoplus_ {i \in \mathrm{I}} \mathrm{E} _ {i} \Bigr) = \operatorname{Card} (\mathrm{I})
\]

(formula (1) of V, p. 377 and cor. 2 of V, p. 387). For every finite subset L of J, we then have

\[
\begin{array}{c} \operatorname{Card} (\mathrm{L}) = \dim \operatorname{Hom} _ {\mathrm{G}} \Bigl (\mathrm{E} _ {\pi}, \bigoplus_ {j \in \mathrm{L}} \mathrm{F} _ {j} \Bigr) \\ \leqslant \dim \operatorname{Hom} _ {\mathrm{G}} (\pi , \operatorname{M} _ {\pi} (\varrho)) = \operatorname{Card} (\mathrm{I}) \end{array}
\]

(loc. cit.). This proves that J is then finite and that \(\text{Card(I)} = \text{Card(J)}\) , as desired. Likewise, if J is finite, then I is finite and has the same cardinal.

Now suppose that I and J are infinite. For \(j \in J\) , denote by \(p_{j}\) the orthogonal projector of E with image \(F_{j}\) . For \(i \in I\) , let \(x_{i}\) be a nonzero element of \(E_{i}\) ; it is a cyclic vector of \(E_{i}\) . Observe that since \(p_{j}\), induced by passage to the subspaces, is a G-morphism from \(E_{i}\) to \(F_{j}\) , the space \(p_{j}(E_{i})\) is zero if and only if \(p_{j}(x_{i}) = 0\) (indeed, the space \(E_{i} \cap \text{Ker}(p_{j})\) is either zero or equal to \(E_{i}\) ).

For every \(j \in J\) , there exists an element \(f(j) \in I\) such that \(p_{j}(\mathrm{E}_{f(j)})\) is not reduced to 0 (otherwise, the orthogonal projector \(p_{j}\) would be zero on \(\mathrm{M}_{\pi}(\varrho)\) ). We have thus defined a map f from J to I. For every \(i \in I\) , the set \(f^{-1}(i)\) is countable. Indeed, this set is contained in the set of \(j \in J\) such that \(p_{j}(\mathrm{E}_{i})\) is nonzero, that is, those such that \(p_{j}(x_{i}) \neq 0\) . Now, by (EVT, V, p. 20, cor. 2), we have

\[
\sum_ {j \in \mathrm{J}} \| p _ {j} (x _ {i}) \| ^ {2} = \| x _ {i} \| ^ {2},
\]

hence the set of \(j \in J\) such that \(p_{j}(x_{i}) \neq 0\) is countable. By E, III, p. 50, prop. 4, we conclude that \(\text{Card}(\text{J}) = \text{Card}(f(\text{J})) \leqslant \text{Card}(\text{I})\) . By interchanging the roles of I and J, we conclude that \(\text{Card}(\text{I}) = \text{Card}(\text{J})\) .

COROLLARY. --- Let G be a topological group and \(\varrho\) be a unitary representation of G in a complex Hilbert space E. The Hilbert sum of the \(\pi\) -isotypic components of G for \(\pi \in \widehat{\mathbf{G}}\) is the largest semisimple subrepresentation of \(\varrho\) .

Indeed, the \(\pi\) -isotypic components of G for \(\pi \in \widehat{\mathbf{G}}\) are pairwise orthogonal (prop. 8, \(a\) ), and each is semisimple (prop. 10, \(a\) ), hence the Hilbert sum F of the spaces \(\mathrm{M}_{\pi}(\varrho)\) for \(\pi \in \widehat{\mathbf{G}}\) defines a semisimple subrepresentation of \(\varrho\) . Since moreover every irreducible subrepresentation of \(\varrho\) is a subrepresentation of an isotypic component of \(\varrho\) , the corollary follows.

DEFINITION 10. --- Let G be a topological group. Let \(\varrho\) be a unitary representation of G in a complex Hilbert space E and \(\pi\) be an irreducible unitary representation of G in a complex Hilbert space.

One calls the multiplicity of \(\pi\) in \(\varrho\) the cardinal of the set I for any family \((\mathrm{E}_i)_{i\in \mathrm{I}}\) of closed subspaces of E stable under G such that the subrepresentation of \(\varrho\) induced in \(\mathrm{E}_i\) is isomorphic to \(\pi\) for every \(i\in I\) and such that \(\mathrm{M}_{\pi}(\varrho)\) is the Hilbert sum of the subspaces \(\mathrm{E}_i\) .

If the multiplicity of \(\pi\) in \(\varrho\) is finite, then it is equal to the dimension of the space \(\mathrm{Hom}_{\mathrm{G}}(\pi, \varrho)\) (resp. to the dimension of \(\mathrm{Hom}_{\mathrm{G}}(\varrho, \pi)\) ) by formula (1) of V, p. 377 and corollary 2 of V, p. 387. This is not always the case in general.

Remark. --- It is possible that a unitary representation \(\varrho\) is nonzero but that all isotypic components of \(\varrho\) relative to all irreducible representations of G are zero (cf. V, p. 426, remark).

\section{REPRESENTATIONS OF LOCALLY COMPACT GROUPS}

In this paragraph, vector spaces are over the field K = R or C. We denote by G a locally compact group equipped with a left Haar measure \(\mu\) . We denote \(\mathcal{L}^{p}(\mathrm{G}) = \mathcal{L}_{\mathbf{C}}^{p}(\mathrm{G}, \mu)\) and \(\mathrm{L}^{p}(\mathrm{G}) = \mathrm{L}_{\mathbf{C}}^{p}(\mathrm{G}, \mu)\) for all \(p \in [1, +\infty]\) .

\subsection{Continuity of certain representations}

PROPOSITION 1. — Let H be a locally compact group. Let \(\varrho_{1}\) (resp. \(\varrho_{2}\) ) be a continuous representation of G (resp. of H) in a locally convex separated K-vector space E \(_{1}\) (resp. E \(_{2}\) ). Let F denote the space \(\mathcal{L}(\mathrm{E}_{1};\mathrm{E}_{2})\) endowed with the topology of compact convergence. The representation \(\varrho\) of G×H into F defined by \(\varrho(g,h)u=\varrho_{2}(h)\circ u\circ\varrho_{1}(g^{-1})\) for \((g,h)\in\mathrm{G}\times\mathrm{H}\) is continuous.

Denote by \(\mathcal{L}_c(\mathrm{E}_1;\mathrm{E}_2)\) the space \(\mathcal{L}(\mathrm{E}_1;\mathrm{E}_2)\) equipped with the topology of compact convergence. Let \(\mathfrak{F}_1\) (resp. \(\mathfrak{F}_2,\mathfrak{F}_3\) ) be a filter in G converging to \(e\) (resp. a filter in \(\mathcal{L}_c(\mathrm{E}_1;\mathrm{E}_2)\) converging to 0, a filter in H converging to \(e\) ). Since G and H are locally compact, there exist elements \(\mathrm{C}\in \mathfrak{F}_1\) and \(\mathrm{D}\in \mathfrak{F}_3\) which are relatively compact. The set \(\varrho_1(\mathrm{C}^{-1})\) is equicontinuous in \(\mathcal{L}(\mathrm{E}_1)\) (cf. INT, VIII, p. 129, § 2, no. 1, remark 2, \(a'\) )). By prop. 9 of EVT, III, p. 33 and prop. 4 of EVT, III, p. 31, the map defined by \((u,v)\mapsto u\circ v\) of \(\mathcal{L}_c(\mathrm{E}_1;\mathrm{E}_2)\times \varrho_1(\mathrm{C}^{-1})\) in \(\mathcal{L}_c(\mathrm{E}_1;\mathrm{E}_2)\) is continuous. The filter base \(\mathfrak{F}_2\circ \varrho_1(\mathfrak{F}_1^{-1})\) therefore converges to 0 in \(\mathcal{L}_c(\mathrm{E}_1;\mathrm{E}_2)\) . The set \(\varrho_2(\mathrm{D})\) is equicontinuous in \(\mathcal{L}(\mathrm{E}_2)\) (cf. INT, VIII, p. 129, § 2, n° 1, rem. 2), and for every \(x\in \mathrm{E}_2\) , the set \(\varrho_2(\mathrm{D})x\subset \mathrm{E}_2\) is relatively compact. Consequently, \(\varrho_2(\mathrm{D})\) is relatively compact in \(\mathcal{L}(\mathrm{E}_2)\) endowed with the topology of compact convergence (TG, X, p. 18, cor. 1). The map defined by \((u,v)\mapsto u\circ v\) of \(\overline{\varrho_2(\mathrm{D})}\times \mathcal{L}_c(\mathrm{E}_1;\mathrm{E}_2)\) in \(\mathcal{L}_c(\mathrm{E}_1;\mathrm{E}_2)\) is continuous by prop. 9 of EVT, III, p. 33 and prop. 4 of EVT, III, p. 31. Hence the filter base \(\varrho (\mathfrak{F}_1\times \mathfrak{F}_3)(\mathfrak{F}_2)\) converges to 0 in \(\mathcal{L}_c(\mathrm{E}_1;\mathrm{E}_2)\) . This implies the assertion.

COROLLARY. --- Let \(\varrho\) be a continuous representation of G on a separated locally convex K-vector space E. The contragredient representation \(\breve{\varrho}\) is continuous when E' is endowed with the topology of compact convergence.

This follows from the proposition by taking \(H = \{e\}\) , \(\varrho_{1} = \varrho\) and \(\varrho_{2}\) the trivial representation on K.

Remark. --- The contragredient representation is not necessarily continuous when E' is endowed with the strong topology (cf. INT, VIII, p. 191, § 2, exercise 3, d)). One can show that it is continuous if the space E is semi-reflexive (loc. cit., c)).

\subsection{Extension of representations to spaces of measures}

In this section, we assume that K = C.

Let \(\varrho\) be a continuous linear representation of G in a separated locally convex quasi-complete space E. Let us denote by \(\mathcal{M}_{c}(\mathrm{G})\) the space of measures with compact support on G equipped with the topology of compact convergence; it is the dual of the space \(\mathcal{C}(\mathrm{G})\) . For every measure \(\nu \in \mathcal{M}_{c}(\mathrm{G})\) we set

\[
\varrho (\nu) = \int_ {\mathrm{G}} \varrho (g) d \nu (g).
\]

This is an element of \(\mathcal{L}(\mathrm{E})\) . In particular, we have \(\varrho(\varepsilon_g) = \varrho(g)\) for all \(g \in \mathrm{G}\) .

According to INT, VIII, p. 136, § 2, no. 6, the map \(\nu \mapsto \varrho(\nu)\) is a continuous linear map from \(\mathcal{M}_c(\mathrm{G})\) to the space \(\mathcal{L}(\mathrm{E})\) equipped with the topology of compact convergence. According to INT, VIII, p. 145, § 3, no. 3, prop. 11, it is a morphism of unital algebras.

Let \(x \in \mathrm{E}\) . The map from \(\mathcal{M}_c(\mathrm{G})\) to E defined by \(\nu \mapsto \varrho(\nu)x\) is continuous when \(\mathcal{M}_c(\mathrm{G})\) is equipped with the topology of compact convergence (INT, VI, p. 27, § 1, no. 7, prop. 16 and EVT, III, p. 31, prop. 4).

For \(f \in \mathcal{L}^{1}(\mathrm{G})\) vanishing outside a compact subset of G, the measure \(f \cdot \mu\) has compact support and we shall denote by \(\varrho^{\mu}(f)\) , or \(\varrho(f)\) when no ambiguity is possible, the endomorphism \(\varrho(f \cdot \mu)\) of E. This endomorphism depends only on the class \(\widetilde{f}\) of \(f\) in \(\mathrm{L}^{1}(\mathrm{G})\) , and we shall also denote it by \(\varrho^{\mu}(\widetilde{f})\) or \(\varrho(\widetilde{f})\) .

Similarly, let \(\varrho\) be a continuous bounded linear representation of G in a Banach space E. For every bounded measure \(\nu \in \mathcal{M}^{1}(\mathrm{G})\) we set

\[
\varrho (\nu) = \int_ {\mathrm{G}} \varrho (g) d \nu (g).
\]

According to INT, VIII, loc. cit., the map \(\nu \mapsto \varrho(\nu)\) is a continuous unital morphism of Banach algebras from the algebra \(\mathcal{M}^{1}(G)\) of bounded complex measures on G into the Banach algebra \(\mathcal{L}(E)\) .

Let \(\rho\) be the function \(g\mapsto\|\varrho(g)\|\) on G; the algebra denoted \(\mathcal{M}^{\rho}\) in INT, VIII, p. 145, prop. 11 (whose elements are the measures \(\nu\) such that \(\rho\in\mathcal{L}^{1}(\nu)\) ) coincides with the Banach algebra \(\mathcal{M}^{1}(G)\) . Indeed, set \(M = \sup \rho\) . We have \(M > 0\) since \(\rho(e)=1\) , and \(M^{-1} \leqslant \rho \leqslant M\) since \(\|\varrho(e)\|\leqslant\|\varrho(g)\|\|\varrho(g^{-1})\|\) for all \(g\in G\) ; thus \(\rho\in\mathcal{L}^{1}(\nu)\) if and only if \(\nu\) is a bounded measure.

If \(f \in \mathcal{L}^{1}(\mathrm{G})\) , one also denotes \(\varrho^{\mu}(f)\) or simply \(\varrho(f)\) the endomorphism \(\varrho(f \cdot \mu)\) , and likewise for the class \(\tilde{f}\) of \(f\) in \(\mathrm{L}^{1}(\mathrm{G})\) .

Lemma 1. --- Let \(\varrho\) be a unitary representation of G in a Hilbert space E. The map \(\nu \mapsto \varrho(\nu)\) of \(\mathcal{M}^{1}(\mathrm{G})\) into \(\mathcal{L}(\mathrm{E})\) is a unital morphism of involutive Banach algebras.

By what precedes, it suffices to show that the morphism \(\nu \mapsto \varrho(\nu)\) is involutive. Let \(\nu \in \mathcal{M}^{1}(\mathrm{G})\) . The measure \(\nu^{*}\) is the conjugate measure of the measure \(\check{\nu}\) (I, p. 99, example 4). According to the definition of the conjugate measure (INT, III, p. 52, § 1, no. 5), one computes

\[
\begin{array}{c} \langle x \mid \varrho (\nu) y \rangle = \int_ {\mathrm{G}} \langle x \mid \varrho (g) y \rangle d \nu (g) \\ = \int_ {\mathrm{G}} \langle \varrho (g ^ {- 1}) x \mid y \rangle d \nu (g) = \int_ {\mathrm{G}} \langle \varrho (g) x \mid y \rangle d \check {\nu} (g) \\ = \overline {{\int_ {\mathrm{G}} \langle y \mid \varrho (g) x \rangle d \nu^ {*} (g)}} = \langle \varrho (\nu^ {*}) x \mid y \rangle \end{array}
\]

for all \(x\) and \(y\) in E, whence \(\varrho(\nu)^{*} = \varrho(\nu^{*})\) .

Let \(\varrho\) be a continuous linear representation (resp. continuous and bounded) of G in a locally convex quasi-complete space E (resp. a Banach space E). If F is a closed subspace of E defining a subrepresentation \(\varrho_{\mathrm{F}}\) of \(\varrho\) , then for every measure \(\nu \in \mathcal{M}_c(\mathrm{G})\) (resp. \(\nu \in \mathcal{M}^1 (\mathrm{G})\) ) the linear map \(\varrho (\nu)\) leaves the subspace F stable and coincides with \(\varrho_{\mathrm{F}}(\nu)\) on F.

Conversely, let \(F \subset E\) be a closed subspace, stable under the linear maps \(\varrho(f)\) for every function \(f \in \mathcal{H}(G)\) . Then F defines a subrepresentation of \(\varrho\) (INT, VIII, p. 139, § 2, no. 7, prop. 10).

Recall (A, VIII, p. 388) that a function \(f \colon \mathbf{G} \to \mathbf{C}\) is said to be central if, for all \(g\) and \(h\) in \(\mathbf{G}\) , we have \(f(gh) = f(hg)\) . This amounts to saying that \(f\) is invariant under conjugation.

DEFINITION 1. --- A bounded measure \(\nu \in \mathcal{M}^{1}(\mathrm{G})\) is said to be central if one has \(\varepsilon_{g} * \nu = \nu * \varepsilon_{g}\) for all \(g \in \mathrm{G}\) .

If G is unimodular and \(f \in \mathcal{C}(\mathrm{G})\) is \(\mu\) -integrable, the measure \(f \cdot \mu\) is central if and only if f is central.

Let \(\varrho\) be a continuous bounded representation of G in a Banach space E (resp. a continuous representation of G in a quasi-complete separated locally convex space E). For every bounded central measure \(\nu\) on G (resp. any compactly supported central measure \(\nu\) on G), the linear map \(\varrho(\nu)\) is a G-morphism from \(\varrho\) to \(\varrho\) . Indeed, for every \(g \in G\) , we have

\[
\varrho (\nu) \varrho (g) = \varrho (\nu * \varepsilon_ {g}) = \varrho (\varepsilon_ {g} * \nu) = \varrho (g) \varrho (\nu).
\]

\subsection{Criterion for semisimplicity}

PROPOSITION 2. --- Let \(\varrho\) be a unitary representation of G in a complex Hilbert space E. Let \(\mathfrak{F}\) be a filter base on \(\mathcal{M}_{c}(\mathrm{G})\) converging to the measure \(\varepsilon_{e}\) for the topology of compact convergence. Suppose that there exists \(\mathrm{M} \in \mathfrak{F}\) such that \(\varrho(\nu)\) is a compact endomorphism of E for every \(\nu \in \mathrm{M}\) .

The unitary representation \(\varrho\) is then semisimple and every irreducible unitary representation of G has finite multiplicity in \(\varrho\) .

Let us first prove a lemma.

Lemma 2. --- Let \(\varrho_{1}\) be a nonzero unitary representation of G isomorphic to a subrepresentation of \(\varrho\) . There exists a measure \(\nu \in M\) such that \(\varrho_{1}(\check{\nu} * \nu)\) is a nonzero compact Hermitian endomorphism of the space of \(\varrho_{1}\) . In particular, there exists a nonzero real number \(\lambda\) such that the eigenspace of \(\varrho_{1}(\check{\nu} * \nu)\) relative to \(\lambda\) is nonzero.

We may assume that \(\varrho_{1}\) is a subrepresentation of \(\varrho\) . Let E \(_{1}\) be its space. For every measure \(\nu \in \mathrm{M}\) , the endomorphism \(\varrho_{1}(\nu)\) is deduced from \(\varrho(\nu)\) by passing to subspaces. The hypothesis and prop. 3 of III, p. 5 therefore imply that \(\varrho_{1}(\nu)\) is compact.

Since \(\mathfrak{F}\) converges to the measure \(\varepsilon_{e}\) in \(\mathcal{M}_{c}(\mathrm{G})\) endowed with the topology of compact convergence and the space \(\mathrm{E}_{1}\) is nonzero, there exists \(\nu\in\mathrm{M}\) such that \(v=\varrho_{1}(\nu)\) is a nonzero endomorphism of \(\mathrm{E}_{1}\) (cf. n°2 of V, p. 400). The endomorphism \(u=\varrho_{1}(\check{\nu}*\nu)\) is equal to \(v^{*}\circ v\) (lemma 1 of V, p. 401); it is therefore nonzero (EVT, V, p. 39, prop. 2), Hermitian and compact, since \(v\) is compact.

Since \(u\) is Hermitian and nonzero, its spectrum is not reduced to zero (example 1 of I, p. 110). Finally, since \(u\) is compact, every \(\lambda \in \mathrm{Sp}(u)\setminus\{0\}\) is an eigenvalue of \(u\) (prop. 2 of III, p. 83).

Let us now prove the proposition.

We will use prop. 7 of V, p. 391 to establish that \(\varrho\) is semisimple. Let \(\varrho_{1}\) be a nonzero subrepresentation of \(\varrho\) and \(\mathrm{E}_{1} \subset \mathrm{E}\) its space; we must show that \(\varrho_{1}\) contains an irreducible subrepresentation.

Let \(\nu \in \mathrm{M}\) such that \(u = \varrho_{1}(\check{\nu}*\nu)\) is a nonzero compact Hermitian endomorphism of the space of \(\varrho_{1}\) and let \(\lambda\) be nonzero such that the eigenspace F of \(u\) relative to \(\lambda\) is nonzero (lemma 2). The space F is finite-dimensional (prop. 5 of III, p. 90). There exists a subrepresentation \(\varrho_{2}\) of \(\varrho_{1}\) on a subspace \(\mathrm{E}_{2}\) of \(\mathrm{E}_{1}\) such that \(\mathrm{E}_{2} \cap \mathrm{F}\) is nonzero and of minimal dimension. Let \(x\) be a nonzero element of \(\mathrm{E}_{2} \cap \mathrm{F}\) and let \(\mathrm{E}_{3}\) be the subrepresentation of \(\varrho_{1}\) generated by \(x\) . We then have \(\mathrm{E}_{3} \subset \mathrm{E}_{2}\) , whence \(\mathrm{E}_{3} \cap \mathrm{F} = \mathrm{E}_{2} \cap \mathrm{F}\) by minimality of the dimension of \(\mathrm{E}_{2} \cap \mathrm{F}\) .

Let us show that the representation \(E_{3}\) is irreducible. Let \(E_{4} \subset E_{3}\) be a closed subspace stable under G. We have \(E_{2} \cap F = (E_{4} \cap F) \oplus (E_{4}^{\circ} \cap F)\) , where \(E_{4}^{\circ}\) denotes the orthogonal complement of \(E_{4}\) in \(E_{3}\) (indeed, if \(y \in E_{2} \cap F\) , let \(y_{4} \in E_{4}\) be its orthogonal projection onto \(E_{4}\) ; since \(E_{4}\) and \(E_{4}^{\circ}\) are stable under u, the vector \(u(y_{4})\) is the orthogonal projection of \(u(y) = \lambda y\) onto \(E_{4}\) , whence \(u(y_{4}) = \lambda y_{4}\) , which means that \(y_{4} \in E_{4} \cap F\) , and the assertion follows). The minimality of the dimension of \(E_{2} \cap F\) then implies that either \(E_{4} \cap F = E_{2} \cap F\) , or \(E_{4}^{\circ} \cap F = E_{2} \cap F\) . In the first case, we have \(x \in E_{4}\) , whence \(E_{4} = E_{3}\) , and in the second, we have \(x \in E_{4}^{\circ}\) , whence \(E_{4}^{\circ} = E_{3}\) and \(E_{4} = \{0\}\) . The representation \(E_{3}\) is therefore irreducible.

We conclude by prop. 7 of V, p. 391 that the representation \(\varrho\) is semisimple.

Let \(\pi\) be an irreducible unitary representation of G whose multiplicity in \(\varrho\) is nonzero; let \(E_{\pi}\) be its space. There exists a measure \(\nu \in M\) and a nonzero real number \(\lambda\) such that \(u = \pi (\check{\nu} * \nu)\) is a nonzero compact Hermitian endomorphism of \(E_{\pi}\) and the eigenspace of \(u\) relative to \(\lambda\) is nonzero (lemma 2). Let F be the eigenspace of \(v = \varrho (\check{\nu} * \nu)\) relative to \(\lambda\) . For every subrepresentation \(E_1\) of E isomorphic to \(\pi\) , the endomorphism of \(E_1\) deduced from \(v\) by passing to subspaces is identified with \(u\) , and the intersection of \(E_1\) and F is nonzero. Thus, the multiplicity of \(\pi\) in \(\varrho\) is less than the dimension of the space F, which is finite (prop. 5 of III, p. 90).

Example. --- Suppose that G is unimodular, for example that G is a real semisimple Lie group. Let \(\Gamma \subset G\) be a discrete subgroup such that the quotient \(X = \Gamma \backslash G\) is compact. The group G acts on X on the right by multiplication. We denote by \(\beta\) the counting measure on \(\Gamma\) and we set \(\widetilde{\mu} = \mu / \beta\) (INT, VII, p. 44, § 2, n° 2, def. 1); this is a bounded G-invariant measure on X.

For every \(f \in \mathcal{L}^2(\mathrm{X},\widetilde{\mu})\) and every \(g \in \mathrm{G}\) , we define the function \(\varrho(g)f \in \mathcal{L}^2(\mathrm{X},\widetilde{\mu})\) by \(\varrho(g)f(x) = f(x \cdot g)\) . The map \(\varrho(g)\) is a continuous linear map, which induces by passing to quotients a unitary map on \(\mathrm{L}^2(\mathrm{X},\widetilde{\mu})\) , still denoted \(\varrho(g)\) . The map \(\varrho\) is a unitary representation of G in \(\mathrm{L}^2(\mathrm{X},\widetilde{\mu})\) (INT, VII, p. 135, § 2, n° 5, prop. 8).

Let \(\varphi \in \mathcal{K}(\mathrm{G})\) . For all compact subsets \(\mathrm{L}_1\) and \(\mathrm{L}_2\) of \(\mathrm{G}\) , the intersection \(\mathrm{T}\) of \(\Gamma\) and of the compact set \(\mathrm{L}_1\mathrm{Supp}(\varphi)\mathrm{L}_2^{-1}\) is finite. For every \((g,h) \in \mathrm{L}_1 \times \mathrm{L}_2\) , the series \(\sum_{\gamma \in \Gamma} \varphi(g^{-1}\gamma h)\) coincides with the finite sum \(\sum_{\gamma \in \mathrm{T}} \varphi(g^{-1}\gamma h)\) . Since \(\mathrm{G}\) is locally compact, the sum of this series, denoted \(k_{\varphi}(g,h)\) , is a continuous function on \(\mathrm{G} \times \mathrm{G}\) .

Let \((g,h) \in G \times G\) and let \((\gamma, \eta) \in \Gamma \times \Gamma\) . We have \(k_{\varphi}(\gamma g, \eta h) = k_{\varphi}(g, h)\) , hence \(k_{\varphi}\) defines by passing to the quotient a continuous function on \(X \times X\) , which we denote \(\widetilde{k}_{\varphi}\) . We also have \(\widetilde{k}_{\varphi} \in \mathcal{L}^{2}(X \times X, \widetilde{\mu} \otimes \widetilde{\mu})\) , since the space X is assumed compact.

We denote by \(\dot{x}\) the image in X of an element \(x\) of G under the canonical projection. Let \(\varphi \in \mathcal{K}(\mathrm{G})\) . For \(f \in \mathcal{L}^2(\mathrm{X},\widetilde{\mu})\) and \(x \in \mathrm{G}\) , we have

\[
\begin{array}{c} \varrho (\varphi) f (\dot {x}) = \int_ {\mathrm{G}} \varphi (g)   (\varrho (g) f) (\dot {x}) d \mu (g) = \int_ {\mathrm{G}} \varphi (g) f (\dot {x} \cdot g) d \mu (g) \\ = \int_ {\mathrm{G}} \varphi (x ^ {- 1} y) f (\dot {y}) d \mu (y) = \int_ {\Gamma \backslash \mathrm{G}} \Bigl (\sum_ {\gamma \in \Gamma} \varphi (x ^ {- 1} \gamma y) \Bigr) f (\dot {y}) d \widetilde {\mu} (\dot {y}) \end{array}
\]

(INT, VII, p. 46, § 2, no. 3, prop. 5). Since \(\widetilde{k}_{\varphi}\) belongs to the space \(\mathcal{L}^2 (\mathrm{X}\times \mathrm{X},\widetilde{\mu}\otimes \widetilde{\mu})\) , it follows that the endomorphism \(\varrho (\varphi)\) of \(\mathrm{L}^2 (\mathrm{X},\widetilde{\mu})\) coincides with the Hilbert-Schmidt endomorphism with kernel \(\widetilde{k}_{\varphi}\) ; this endomorphism is therefore compact (cor. 1 of III, p. 33).

There exists a sequence \((\varphi_{n})_{n\in\mathbf{N}}\) of functions in \(\mathcal{K}_{+}(\mathrm{G})\) with integral 1 such that \(\varphi_{n}\cdot\mu\) converges to \(\varepsilon_{e}\) in \(\mathcal{M}_{c}(\mathrm{G})\) endowed with the topology of compact convergence (INT, VIII, p. 139, § 2, no. 7, cor. 2). Prop. 2 consequently implies that the representation \(\varrho\) is semisimple and that the multiplicities of the irreducible unitary representations of G in \(\varrho\) are finite.

The irreducible unitary representations of G whose multiplicity in \(\varrho\) is nonzero are called \(\Gamma\) -automorphic representations of the group G.

\subsection{Regular representations}

Let \(p\) be a real number \(\geqslant 1\) and \(\mu'\) a right Haar measure on \(G\) . For every \(g \in G\) and every function \(f \in \mathcal{L}^p(G, \mu)\) , we denote by \(\gamma_G^{(p)}(g)f\) the function \(x \mapsto f(g^{-1}x)\) on \(G\) . The map \(g \mapsto \gamma_G^{(p)}(g)\) is a continuous linear representation of \(G\) in \(\mathcal{L}^p(G)\) . By passing to quotients, it induces a continuous isometric linear representation of \(G\) in \(L^p(G)\) (INT, VIII, p. 135, § 2, no. 5, prop. 8), denoted \(\gamma_G^{(p)}\) .

Similarly, for every \(g \in G\) and every function \(f \in \mathcal{L}^{p}(G, \mu')\) , we denote by \(\boldsymbol{\delta}_{\mathrm{G}}^{(p)}(g)f\) the function \(x \mapsto f(xg)\) on G. The map \(g \mapsto \boldsymbol{\delta}_{\mathrm{G}}^{(p)}(g)\) is a continuous linear representation of G in \(\mathcal{L}^{p}(G, \mu')\) . By passing to quotients, it induces a continuous isometric linear representation of G in \(\mathrm{L}^{p}(\mathrm{G}, \mu')\) (cf. INT, VIII, p. 136, § 2, no. 5).

We say that \(\boldsymbol{\gamma}_{\mathrm{G}}^{(p)}\) is the left regular representation of G in \(\mathcal{L}^{p}(\mathrm{G})\) or \(\mathrm{L}^{p}(\mathrm{G})\) and that \(\boldsymbol{\delta}_{\mathrm{G}}^{(p)}\) is the right regular representation of G in \(\mathcal{L}^{p}(\mathrm{G},\mu^{\prime})\) or \(\mathrm{L}^{p}(\mathrm{G},\mu^{\prime})\) .

Lemma 3. — Let p be a real number ≥ 1. The left (resp. right) regular representation of G in L \(^{p}\) (G,μ) (resp. in L \(^{p}\) (G,μ')) is faithful.

More precisely, let q be the conjugate exponent of p and let g be an element of G such that \(g \neq e\) .

\begin{enumerate}
\def\labelenumi{\alph{enumi})}
\item
  There exists a function \(\varphi\in\mathcal{K}(\mathrm{G})\) , positive and nonzero in \(\mathrm{L}^{q}(\mathrm{G},\mu)\) , such that \(\langle\varphi,\boldsymbol{\gamma}_{\mathrm{G}}^{(p)}(g)\varphi\rangle=0\) ;
\item
  There exists a function \(\varphi\in\mathcal{K}(\mathrm{G})\) , positive and nonzero in \(\mathrm{L}^{q}(\mathrm{G},\mu')\) , such that \(\langle\varphi,\boldsymbol{\delta}_{\mathrm{G}}^{(p)}(g)\varphi\rangle=0\) .
\end{enumerate}

Consider the case of the left regular representation \(\boldsymbol{\gamma}_{\mathrm{G}}^{(p)}\) , that of the right regular representation being similar. The assertion \(a\) ) implies that \(\boldsymbol{\gamma}_{\mathrm{G}}^{(p)}\) is faithful, for if \(g \neq e\) is an element of G, and if \(\varphi\) is as in \(a\) ), one has \(\varphi \neq \boldsymbol{\gamma}_{\mathrm{G}}^{(p)}(g)\varphi\) since \(\langle \varphi, \varphi \rangle = \int_{\mathrm{G}} \varphi^2 > 0\) .

Let us therefore prove \(a\) ). Let \(g \neq e\) in G. Let C be a compact symmetric neighborhood of \(e\) such that \(g \notin \mathrm{C}^2\) and let \(\varphi \in \mathcal{K}(\mathrm{G})\) be a continuous positive function with integral 1 and support contained in C; the function \(\varphi\) is nonzero in \(\mathrm{L}^p(\mathrm{G}, \mu)\) . Since \(\mathrm{C} \cap g^{-1}\mathrm{C} = \emptyset\) , we have

\[
\langle \varphi , \boldsymbol {\gamma} _ {\mathrm{G}} ^ {(p)} (g) \varphi \rangle = \int_ {\mathrm{G}} \varphi (x) \varphi (g ^ {- 1} x) d \mu (x) = 0.
\]

Let \(\varrho\) be a continuous and bounded representation of G in a Banach space E. For every \(f\in \mathrm{L}^1 (\mathrm{G})\) (resp. \(f^{\prime}\in \mathrm{L}^{1}(\mathrm{G},\mu^{\prime}))\) and every \(g\in \mathrm{G}\) , we have

\[
\varrho (g) \varrho (f \cdot \mu) = \varrho (\varepsilon_ {g} * (f \cdot \mu)) = \varrho (\boldsymbol {\gamma} _ {\mathrm{G}} ^ {(1)} (g) f \cdot \mu),\tag{1}
\]

\[
\varrho (f ^ {\prime} \cdot \mu^ {\prime}) \varrho (g) = \varrho ((f ^ {\prime} \cdot \mu^ {\prime}) * \varepsilon_ {g}) = \varrho (\boldsymbol {\delta} _ {\mathrm{G}} ^ {(1)} (g ^ {- 1}) f ^ {\prime} \cdot \mu ^ {\prime}),\tag{2}
\]

(INT, VIII, p. 144, § 3, no. 2, formula (5)).

Suppose G is unimodular, and set \(\mu' = \mu\) . The biregular representation of G in \(\mathcal{L}^{p}(\mathrm{G})\) (resp. in \(\mathrm{L}^{p}(\mathrm{G})\) ) is the representation \(\varrho_{\mathrm{G}}^{(p)}\) of \(G \times G\) in \(\mathcal{L}^{p}(\mathrm{G})\) (resp. in \(\mathrm{L}^{p}(\mathrm{G})\) ) such that

\[
\boldsymbol {\varrho} _ {\mathrm{G}} ^ {(p)} (g _ {1}, g _ {2}) = \boldsymbol {\gamma} _ {\mathrm{G}} ^ {(p)} (g _ {1}) \boldsymbol {\delta} _ {\mathrm{G}} ^ {(p)} (g _ {2}) = \boldsymbol {\delta} _ {\mathrm{G}} ^ {(p)} (g _ {2}) \boldsymbol {\gamma} _ {\mathrm{G}} ^ {(p)} (g _ {1}).
\]

It is a continuous linear representation (lemma 1 of V, p. 377). It satisfies

\[
(\boldsymbol {\varrho} _ {\mathrm{G}} ^ {(p)} (g _ {1}, g _ {2}) f) (x) = f (g _ {1} ^ {- 1} x g _ {2})
\]

for all \(f \in \mathcal{L}^p(\mathrm{G})\) , all \((g_1, g_2) \in \mathrm{G} \times \mathrm{G}\) and every \(x \in \mathrm{G}\) .

Note. --- The biregular representation of G in \(\mathrm{L}^{p}(\mathrm{G},\mu)\) is not necessarily faithful; its kernel is the image of the center of G under the map \(g\mapsto(g,g)\) (exercise 4 of V, p. 487).

When \(p = 2\) the left regular representation \(\gamma_{\mathrm{G}}^{(2)}\) of G in the complex Hilbert space \(\mathrm{L}^2 (\mathrm{G},\mu)\) is unitary, since it is isometric. Likewise, the right regular representation \(\delta_{\mathrm{G}}^{(2)}\) in \(\mathrm{L}^2 (\mathrm{G},\mu')\) is unitary.

We will simply denote \(\gamma_{\mathrm{G}} = \gamma_{\mathrm{G}}^{(2)}\) and \(\delta_{\mathrm{G}} = \delta_{\mathrm{G}}^{(2)}\) and these representations will be called the left and right regular representations of G.

If G is unimodular, the biregular representation \(\varrho_{\mathrm{G}}^{(2)}\) of \(G \times G\) in \(\mathrm{L}^{2}(\mathrm{G}, \mu)\) is unitary. We will simply denote it by \(\varrho_{G}\) .

Lemma 4. --- Let \(p\) be a real number \(\geqslant 1\) . For \(f \in \mathcal{L}^1(\mathrm{G})\) , the linear map \(\gamma_{\mathrm{G}}^{(p)}(f)\) coincides with the endomorphism of \(\mathrm{L}^p(\mathrm{G})\) defined by \(\varphi \mapsto f *^\mu \varphi\) .

This follows from INT, VIII, p. 157, § 4, n°2, prop. 6, taking into account formula (14) of INT, VIII, p. 165.

Note. --- Recall that, for any function \(f\) on G, one defines the function \(\check{f}\) on G by \(\check{f}(g) = f(g^{-1})\) for all \(g \in \mathrm{G}\) (INT, VII, p. 12, § 1, n° 1, formula (12)).

One checks that if \(f \in \mathcal{L}^{1}(\mathrm{G})\) , the linear map \(u = \boldsymbol{\delta}_{\mathrm{G}}^{(p)}(f)\) satisfies the relation

\[
\widetilde {u (\varphi)} = f * \check {\varphi}
\]

for all \(\varphi\in\mathrm{L}^{p}(\mathrm{G},\mu^{\prime})\) .